% !TEX root = ../main.tex
% Structural Analysis. Two structural results on swarm-UAV AoI under a shared,
% finite-source (machine-repair) charging queue.

We now establish the two structural properties that separate the swarm regime
from the single-UAV base case: the peak AoI is non-monotone in the swarm size
(Proposition~\ref{prop:ushape}), and the AoI-optimal station placement inverts
from coverage-driven to traffic-driven once the swarm contends for ports
(Proposition~\ref{prop:inversion}). Throughout, we work with the per-UAV peak
AoI of the homogeneous single-station cell derived in Section~\ref{sec:model},
\begin{equation}
\label{eq:aoi-decomp}
\Aoi(M) \;=\; \underbrace{t_{\mathrm{loop}}(M)}_{c_1(M)}
\;+\; \underbrace{\Wq(M,c)}_{\text{contention}}
\;+\; \underbrace{\tfrac{2\,\mathrm{dist}}{V} + \tau_{\mathrm{charge}}}_{c_0},
\end{equation}
where $M$ UAVs share one field and one station of $c$ ports, $t_{\mathrm{loop}}(M)$
is the patrol (revisit) time per UAV, $\Wq(M,c)$ is the mean charging wait of the
finite-source queue $M/M/c//M$ of Section~\ref{sec:model} (population $M$,
per-source operating rate $\lambda=1/T_{\mathrm{up}}$, port rate
$\mu=1/\tau_{\mathrm{charge}}$), and $c_0$ collects the recharge round trip and the
charging time, which do not depend on $M$.

\subsection{Congestion threshold: AoI is non-monotone in swarm size}

We first isolate the two competing forces in \eqref{eq:aoi-decomp}: the coverage
term $c_1(M)$ that a larger swarm improves, and the wait $\Wq(M,c)$ that a larger
swarm degrades.

\begin{assumption}[Coverage scaling]
\label{ass:coverage}
The patrol term $c_1(M)=t_{\mathrm{loop}}(M)$ is positive, strictly decreasing,
and convex in $M$, with $c_1(M)=\Theta(1/M)$.
\end{assumption}

\noindent
The scaling is the standard sub-tour argument: $M$ UAVs partition $K$ sensors of a
field of area $\mathcal{A}$ into $M$ sub-fields of area $\mathcal{A}/M$ and
$K/M$ sensors each. The length of a near-optimal tour through $n$ points uniformly
spread over an area $\mathcal{A}'$ scales as $\Theta(\sqrt{\mathcal{A}'\,n})$, so
each sub-tour has length $\Theta(\sqrt{(\mathcal{A}/M)(K/M)})=\Theta(\sqrt{\mathcal{A}K}/M)$,
giving $t_{\mathrm{loop}}(M)=\Theta(1/M)$. Convexity and monotonicity of $1/M$
carry over to the leading term, and we take them as the modeling assumption above.

\begin{lemma}[Monotone, saturating charging wait]
\label{lem:wq}
In the finite-source queue $M/M/c//M$, the mean wait $\Wq(M,c)$ is non-decreasing
in the population $M$, is strictly increasing once $M>c$, and satisfies the bound
\begin{equation}
\label{eq:wq-bound}
0 \;\le\; \Wq(M,c) \;\le\; \frac{M-c}{c\,\mu}, \qquad M>c .
\end{equation}
\end{lemma}

\begin{IEEEproof}
Let $\{n_t\}$ denote the birth-death chain of Section~\ref{sec:model} on states
$n\in\{0,\dots,M\}$ (UAVs at the station), with up-rates $(M-n)\lambda$ and
down-rates $\min(n,c)\mu$ \cite{grossharris}. Increasing the population from $M$
to $M+1$ adds one source: at every state the up-rate weakly increases (from
$(M-n)\lambda$ to $(M{+}1-n)\lambda$) while the down-rate is unchanged. A standard
monotone coupling of the two chains keeps the $(M{+}1)$-chain pathwise at least as
high as the $M$-chain, so the stationary number in system is stochastically
increasing in $M$.

We convert this into a statement about the wait through the arriving-UAV view,
rather than through the queue length alone: because the effective throughput
$\lambda_{\mathrm{eff}}(M)$ also grows with $M$, monotonicity of the ratio
$\Wq=L_{\mathrm q}/\lambda_{\mathrm{eff}}$ does not follow from monotonicity of
$L_{\mathrm q}$ by itself. By the arrival theorem for finite-source
(machine-repair) queues \cite{grossharris}, a UAV that requests a charge in the
$M$-population system finds the station distributed as the \emph{stationary}
distribution of the $(M{-}1)$-population system. A UAV that finds $n$ present is
served after $(n-c+1)^{+}$ successive departures, each exponential of rate $c\mu$
while the $c$ servers stay busy, so its conditional mean wait is
$(n-c+1)^{+}/(c\mu)$, which is non-decreasing in $n$. Averaging a non-decreasing
function against a distribution that is stochastically increasing in $M$,
$\Wq(M,c)$ is non-decreasing in $M$ (and equals the Little's-law form
$L_{\mathrm q}/\lambda_{\mathrm{eff}}$ evaluated in Section~\ref{sec:model}). The
increase is strict once $M>c$, because states $n>c$ then carry positive
probability in the arriving-UAV distribution and contribute a strictly positive
wait. For \eqref{eq:wq-bound}, an arriving UAV finds at most $M-1$ others present,
so $n\le M-1$ and its conditional wait is at most
$(M-1-c+1)/(c\mu)=(M-c)/(c\mu)$; taking the expectation gives
$\Wq(M,c)\le (M-c)/(c\mu)$.
\end{IEEEproof}

\noindent
Two features of Lemma~\ref{lem:wq} drive the result. Below $c$ the wait is
essentially zero (no UAV ever queues), and above $c$ it climbs, ultimately at the
linear rate $\Theta(M/(c\mu))$ of \eqref{eq:wq-bound}: its forward differences
$\Delta\Wq(M)\triangleq\Wq(M{+}1,c)-\Wq(M,c)$ are non-negative and rise from
$\Delta\Wq\approx0$ in the underloaded region toward the saturation slope
$1/(c\mu)$. We take this saturating shape as an explicit assumption.

\begin{assumption}[Saturating charging increments]
\label{ass:wq-convex}
For each fixed $c$, the increment $\Delta\Wq(M,c)$ is non-decreasing in $M$ (so
$\Wq(\cdot,c)$ is discretely convex), and $\Delta\Wq(M,c)$ is non-increasing in
$c$.
\end{assumption}

\noindent
Both parts reflect the same mechanism: added contention (larger $M$) steepens the
wait, while added ports (larger $c$) flatten it and defer saturation. They are
verified numerically across the tested field and capacity configurations in
Section~\ref{sec:exp}.

\begin{proposition}[Congestion threshold and provisioning law]
\label{prop:ushape}
Under Assumption~\ref{ass:coverage} and Lemma~\ref{lem:wq}, the peak AoI
$\Aoi(M)$ in \eqref{eq:aoi-decomp} is quasi-convex in the swarm size $M$ and
admits a unique interior minimizer
\begin{equation}
M^\star \;=\; \argmin_{M\in\{1,2,\dots\}} \Aoi(M),
\end{equation}
so that adding UAVs strictly reduces AoI for $M<M^\star$ (coverage dominates) and
strictly increases it for $M>M^\star$ (charging congestion dominates). Moreover
$M^\star$ is non-decreasing in the port capacity $c$: enlarging a station never
lowers the optimal swarm size.
\end{proposition}

\begin{IEEEproof}
Since $c_0$ is constant in \eqref{eq:aoi-decomp}, work with the forward difference
\begin{equation}
\label{eq:delta-aoi}
\Delta\Aoi(M) \;=\; \Delta c_1(M) \;+\; \Delta\Wq(M),
\end{equation}
where $\Delta c_1(M)=c_1(M{+}1)-c_1(M)$. By Assumption~\ref{ass:coverage},
$c_1$ is strictly decreasing and convex, so $\Delta c_1(M)<0$ and $\Delta c_1(M)$
is non-decreasing (it rises toward $0$ as $M\to\infty$). By Lemma~\ref{lem:wq}
$\Delta\Wq(M)\ge0$, and by Assumption~\ref{ass:wq-convex} it is non-decreasing.
Hence the right-hand side of
\eqref{eq:delta-aoi} is a sum of two non-decreasing terms, so $\Delta\Aoi(M)$ is
non-decreasing in $M$: the sequence $\Aoi(M)$ is discretely convex.

It remains to check the endpoints of the sign of $\Delta\Aoi$. For small $M$
(in particular $M\le c$) the queue is empty, $\Delta\Wq(M)=0$, while
$\Delta c_1(M)<0$, so $\Delta\Aoi(M)<0$: AoI is still falling. As $M\to\infty$,
$\Delta c_1(M)=\Theta(1/M)-\Theta(1/(M{+}1))\to0^{-}$ while, by
\eqref{eq:wq-bound}, $\Delta\Wq(M)\to 1/(c\mu)>0$, so $\Delta\Aoi(M)>0$
eventually: AoI is rising. A non-decreasing difference sequence that is negative
for small $M$ and positive for large $M$ changes sign exactly once, from negative
to positive. The last index with $\Delta\Aoi(M)<0$ is the unique interior
minimizer $M^\star$, which establishes the U-shape.

For the provisioning law, fix $M$ and increase $c$. The same monotone coupling,
now applied to the down-rates $\min(n,c)\mu$ (which increase with $c$), makes the
stationary occupancy stochastically smaller, so $\Wq(M,c)$ is decreasing in $c$
(more ports serve the same population); by Assumption~\ref{ass:wq-convex} its
increments $\Delta\Wq(M;c)$ are non-increasing in $c$ as well. Enlarging the
station thus both lowers the wait and defers the onset of the steep saturation
region to a larger population. The coverage difference $\Delta c_1(M)$ is
unchanged. Hence for $c'>c$ we have $\Delta\Aoi(M;c')\le\Delta\Aoi(M;c)$ for all
$M$, so the single sign change of $\Delta\Aoi(\cdot;c')$ occurs at an index at
least as large as that of $\Delta\Aoi(\cdot;c)$, i.e. $M^\star(c')\ge M^\star(c)$.
\end{IEEEproof}

\noindent
Proposition~\ref{prop:ushape} says that congestion is a genuine, capacity-limited
effect rather than an artifact of poor routing: even the balanced single-cell wait
of Lemma~\ref{lem:wq} forces a finite $M^\star$. The provisioning law
$M^\star(c)$ non-decreasing is the design counterpart: to operate a larger swarm
at low AoI, one must add ports, not merely UAVs. Both the U-shape and the rightward
shift of $M^\star$ with capacity are confirmed empirically, under optimal (load
balancing) assignment, in Section~\ref{sec:exp}.

\subsection{Placement inversion under contention}

For a single UAV the wait term in \eqref{eq:aoi-decomp} vanishes ($M=1\le c$), so
$\Aoi=c_1(1)+2\,\mathrm{dist}/V+\tau_{\mathrm{charge}}$ depends on placement only
through the recharge detour $\mathrm{dist}$. The AoI-optimal station therefore sits
where the detour is smallest, i.e. near the trajectory centroid: placement is
coverage-driven. The next result shows this rule is not merely loose but can be
strictly wrong for a swarm.

\begin{proposition}[Placement inversion]
\label{prop:inversion}
There exist instances in which the coverage-optimal (minimum-detour) placement is
strictly AoI-suboptimal for a swarm. Concretely, with a per-station port cap and
two candidate sites, opening only the nearest site is beaten by opening both sites
and splitting the swarm, for every swarm size above a finite threshold $M_0$.
\end{proposition}

\begin{IEEEproof}
Consider two candidate sites and a per-station port cap $C_{\mathrm{port}}=c$
(Section~\ref{sec:model}), so a single station can host at most $c$ ports; adding
capacity beyond $c$ requires opening a second station. Site~$1$ is central with
zero extra detour; site~$2$ is offset, adding a round-trip detour
$\delta=2\Delta/V>0$ for any UAV assigned to it. Two placements are compared under
the same swarm and field, so the coverage term $c_1(M)=t_{\mathrm{loop}}$ is
identical in both:
\begin{align}
\text{(A) coverage-optimal:}\quad
&\Aoi_{\mathrm A}(M) = t_{\mathrm{loop}} + \Wq(M,c) + c_0, \\
\text{(B) split:}\quad
&\Aoi_{\mathrm B}(M) = t_{\mathrm{loop}} + \delta \nonumber\\
&\hphantom{\Aoi_{\mathrm B}(M) ={}} + \Wq\!\left(\tfrac{M}{2},c\right) + c_0 ,
\end{align}
where (A) sends all $M$ UAVs to the near station (its $c$ ports serving population
$M$), and (B) opens both stations with $c$ ports each and assigns $M/2$ UAVs to
each; the peak in (B) is attained by the UAVs on the offset station, which carry
the extra detour $\delta$. Subtracting,
\begin{equation}
\Aoi_{\mathrm A}(M)-\Aoi_{\mathrm B}(M)
= \Big[\Wq(M,c)-\Wq\!\left(\tfrac{M}{2},c\right)\Big] - \delta .
\end{equation}
Two bounds make the bracket diverge. For the near station, the throughput cannot
exceed the service ceiling, $\lambda_{\mathrm{eff}}\le c\mu$; combined with the
finite-source cycle identity
$\lambda_{\mathrm{eff}}=M/(T_{\mathrm{up}}+\Wq+\tau_{\mathrm{charge}})$ (each of
the $M$ UAVs alternates operating, waiting, and charging), this gives the linear
lower bound
\begin{equation}
\Wq(M,c)\;\ge\;\frac{M}{c\mu}-T_{\mathrm{up}}-\tau_{\mathrm{charge}} .
\end{equation}
For the split station, the Lemma upper bound \eqref{eq:wq-bound} gives
$\Wq(M/2,c)\le (M/2-c)/(c\mu)\le M/(2c\mu)$. Subtracting,
\begin{equation}
\Wq(M,c)-\Wq\!\left(\tfrac{M}{2},c\right)
\;\ge\;\frac{M}{2c\mu}-T_{\mathrm{up}}-\tau_{\mathrm{charge}},
\end{equation}
which grows without bound in $M$ (the constants $T_{\mathrm{up}}$ and
$\tau_{\mathrm{charge}}$ do not depend on $M$), whereas $\delta$ is fixed.
Therefore there is a finite threshold
\begin{equation}
M_0 = \min\Big\{ M : \Wq(M,c)-\Wq\!\left(\tfrac{M}{2},c\right) > \delta \Big\}
\end{equation}
such that $\Aoi_{\mathrm A}(M)>\Aoi_{\mathrm B}(M)$ for all $M>M_0$. For such $M$
the minimum-detour placement (A) is strictly AoI-suboptimal, so the AoI-optimal
placement must open the offset site and balance the charging load, even at the cost
of a longer detour.
\end{IEEEproof}

\noindent
Proposition~\ref{prop:inversion} formalises the inversion: as the swarm grows, the
optimal placement stops chasing the shortest detour and starts spreading stations
to balance queueing load, so the objective flips from coverage-driven to
traffic-driven. The single-UAV placement rule of the base model is thus not a
conservative approximation but the wrong optimizer for a contended swarm. This
crossover, including the emergence of the load-balanced placement as the field and
swarm scale, is demonstrated empirically in Section~\ref{sec:exp}.
